%% Preamble for the LiSA CVPR 2027 submission.
%% Based on preamble.tex of the official cvpr-org/author-kit (main branch); only additions below.

%% Official support for 2-column teaser figure
\usepackage{cuted} %< for strip environment
\usepackage{currfile} %< for \currfilebase macro
\usepackage{caption} %< for "\captionof" commands (teaser)

%% Tables and figures
\usepackage{multirow}
\usepackage{tikz}
\usetikzlibrary{positioning,fit,calc,arrows.meta,backgrounds}

%% Text layout (enabled to stay within the page limit without altering the official margins)
\usepackage{microtype}

%% Inline annotations (from the kit). Placeholders for experiments that still have to be run are
%% rendered in red so that they cannot be overlooked; see paper/README.md, "Experiments to run".
\newcommand{\red}[1]{{\color{red}#1}}
\newcommand{\todo}[1]{{\color{red}#1}}
\newcommand{\TODO}[1]{\textbf{\color{red}[TODO: #1]}}
\newcommand{\TBD}{{\color{red}\textbf{TBD}}}
\newcommand{\torun}[1]{{\color{red}\textbf{[TO RUN]} #1}}
%% For the camera-ready / final submission, all \torun and \TBD entries must be resolved.

%% Shorthands
\usepackage{xspace}
\newcommand{\ours}{LiSA\xspace}
\newcommand{\gsthree}{GS$^3$\xspace}
\newcommand{\vect}[1]{\mathbf{#1}}
\DeclareMathOperator{\logit}{logit}
\DeclareMathOperator{\softplus}{softplus}

%% Chinese reading edition. XeLaTeX; fonts stay local to the translation.
\usepackage{fontspec}
\setmainfont[Path=../assets/fonts/,AutoFakeBold=2,AutoFakeSlant=0.15]{NotoSerifCJKsc-Regular.otf}
\setsansfont[Path=../assets/fonts/,AutoFakeBold=2,AutoFakeSlant=0.15]{NotoSerifCJKsc-Regular.otf}
\setmonofont[Path=../assets/fonts/,AutoFakeBold=2]{NotoSerifCJKsc-Regular.otf}
\XeTeXlinebreaklocale "zh"
\XeTeXlinebreakskip=0pt plus 1pt
\setlength{\emergencystretch}{1em}
\renewcommand{\baselinestretch}{1.05}
% The upstream style uses legacy TeX fonts for small headings.
\def\elvbf{\normalfont\bfseries\fontsize{11}{13}\selectfont}
\def\tenbf{\normalfont\bfseries\fontsize{10}{12}\selectfont}
\renewcommand{\figurename}{图}
\renewcommand{\tablename}{表}
\renewcommand{\refname}{参考文献}
\renewcommand{\torun}[1]{{\color{red}\textbf{[待开展]} #1}}
\renewcommand{\TBD}{{\color{red}\textbf{待补}}}
\AtBeginDocument{
\crefname{figure}{图}{图}
\Crefname{figure}{图}{图}
\crefname{table}{表}{表}
\Crefname{table}{表}{表}
\crefname{equation}{式}{式}
\Crefname{equation}{式}{式}
\crefformat{section}{第~#2#1#3~节}
\Crefformat{section}{第~#2#1#3~节}
}
\makeatletter
\patchcmd{\abstract}{Abstract}{摘要}{}{}
\patchcmd{\@maketitle}{Anonymous \confName~submission}{\confName~匿名投稿}{}{}
\patchcmd{\@maketitle}{Paper ID}{论文编号}{}{}
\patchcmd{\maketitlesupplementary}{Supplementary Material}{补充材料}{}{}
\makeatother
